% Generated draft: requires mathematical and source review.
\documentclass[11pt]{article}
\usepackage[margin=25mm]{geometry}
\usepackage{fontspec}
\setmainfont{Latin Modern Roman}
\usepackage{amsmath,amssymb,mathtools,unicode-math}
\setmathfont{Latin Modern Math}
\usepackage{enumitem,needspace,xurl,hyperref}
\setlength{\parindent}{0pt}
\setlength{\parskip}{5pt}
\setlength{\emergencystretch}{3em}
\newcommand{\N}{\mathbb{N}}
\newcommand{\Z}{\mathbb{Z}}
\newcommand{\Q}{\mathbb{Q}}
\newcommand{\R}{\mathbb{R}}
\newcommand{\C}{\mathbb{C}}
\newcommand{\F}{\mathbb{F}}
\newcommand{\A}{\mathbb{A}}
\newcommand{\PP}{\mathbb{P}}
\newcommand{\E}{\mathbb{E}}
\newcommand{\eps}{\varepsilon}
\newcommand{\dd}{\,\mathrm{d}}
\begin{document}
\section*{UnsolvedMath ID: 30004985}
\textbf{Source problem code:} catalogue entry
\subsection{Definitions and mathematical statement}
For finite subsets \(A,B\subset\R^3\), write \(A\cong B\) if there exist an orthogonal matrix \(U\in O(3)\) and a vector \(t\in\R^3\) such that \(B=\{Ua+t:a\in A\}\). Thus congruence allows reflections as well as rotations and translations. For a positive integer \(k\), define the \emph{\(k\)-deck} of \(A\) to be the multiset
\[
D_k(A)=\bigl\{\!\bigl\{[S]_{\cong}:S\subseteq A,\ |S|=k\bigr\}\!\bigr\}.
\]
Each unordered \(k\)-element subset contributes once, and congruent subsets contribute with their multiplicities. The points of \(A\) are distinct; no general-position hypothesis is supplied.

The supplied clean statement asks whether a universal positive integer \(k\), independent of the finite point set and its cardinality, suffices for reconstruction up to congruence from this deck. There is an unresolved boundary qualification: the source does not specify how sets with fewer than \(k\) points are handled, although their \(k\)-decks are empty. With the additional restriction that the compared sets have at least \(k\) points, a precise candidate formulation is
\[
\exists k\in\Z_{\ge1}\ \forall A,B\subset\R^3\text{ finite},\qquad
\bigl(|A|\ge k\ \land\ |B|\ge k\ \land\ D_k(A)=D_k(B)\bigr)
\Longrightarrow A\cong B.
\]
The cardinality restriction in this display is an explicit possible repair, not a verified source hypothesis. The original opening question instead asks how large \(k\) must be for every set of \(n\) points to be reconstructible. If cardinality \(n\) is supplied separately, its corresponding precise property, for \(n\ge k\ge1\), is
\[
\forall A,B\subset\R^3\text{ finite},\qquad
\bigl(|A|=|B|=n\ \land\ D_k(A)=D_k(B)\bigr)
\Longrightarrow A\cong B.
\]
The intended boundary convention, whether cardinality is additional reconstruction data, and whether the target is existence of a constant or the optimal value of \(k\) require source review.
\subsection{Short English statement}
The question asks whether one fixed number of points per card suffices to reconstruct every finite configuration in \(\R^3\), up to Euclidean isometry, from the multiplicities of all its congruence classes of cards. The treatment of configurations smaller than a card and the role of a known total cardinality remain unspecified.
\subsection{Sources}
\begin{itemize}
\item Local UnsolvedMath registry: \texttt{data/unsolvedmath/problems.json}, permanent ID \(30004985\), source code \texttt{catalogue entry}, title ``Reconstruction of Point Sets from \(k\)-Decks.'' The supplied record lists proposed year 2022 and no named proposer.
\item Supplied source citation: ``Graph Theory (2023), Oberwolfach Reports, DOI 10.4171/owr/2022/1, pp. 5--78.'' \url{https://doi.org/10.4171/owr/2022/1}. The background describes Oberwolfach Reports as records of mathematical workshops and their open-problem discussions. The supplied literature list also labels this reference ``Oberwolfach Report 1/2022.'' The differing year labels are retained without reconciliation.
\item Supplied literature reference: ``Pebody, Radcliffe, and Scott, finite planar subsets.'' \url{https://digitalcommons.unl.edu/mathfacpub/142/}.
\item Literature assessment dated exactly \texttt{2026-08-22}: the supplied curation labels the problem ``partially solved'' and reports that a universal constant deck size is known for planar point sets, whereas in dimensions at least three only logarithmic upper bounds and the general constant-size conjecture are known; it describes the requested sharp \(k\) as open. This is supplied dated curation, not an independently verified literature assessment. The background expressly requests independent verification of time-sensitive status and cited claims.
\item The supplied record labels its clean statement \texttt{corrected\_verified} and says it was checked against the official report and preserves the report's hypotheses, notation, and open target. These are supplied verification claims, not verification performed for this draft.
\end{itemize}
\end{document}
